\chapter*{前言：本选集如何炼成}

这本小册子把一座由 \textbf{99 篇格点 QCD 研究论文}组成的工作文献库
（\texttt{papers/} 收藏）提炼为其最重要的论文集合。按任务要求，选编
只依据两个来源：

\begin{enumerate}
  \item \textbf{原论文本身}——其中若干篇本身就是领域里程碑；以及
  \item 它们的\textbf{参考文献列表}——我们对每一篇论文的书目逐一
        挖掘，并在整个语料库内统计了引文频次。
\end{enumerate}

一篇论文入选，凭以下两条理由之一（或兼有）：(i)~它是文献库的
\emph{原论文之一}，且其内容本身就是领域里程碑；或 (ii)~它是引文网络的
\emph{枢纽}——被库内论文反复依赖的奠基之作。下文引用的统计数字是对
这一特定语料库的保守估计；它衡量的是\emph{与本库的相关度}，而非全局
影响力（后者只会更大）。

\section*{领域地图}

38 篇入选论文自然组织为八章，它们也正是现代格点部分子物理的八根承重柱：

\begin{center}
\begin{tabular}{clc}
\toprule
章 & 主题 & 篇数\\
\midrule
1 & QCD 与部分子思想：理论根基 & 3\\
2 & 模拟引擎与信号增强 & 4\\
3 & Yang--Mills 梯度流纲领 & 6\\
4 & 从欧氏格点到光锥：LaMET 革命 & 7\\
5 & 重整化攻坚战 & 7\\
6 & 胶子 PDF：两条主线的交汇 & 5\\
7 & TMD 与全球拟合：与唯象学的对接 & 3\\
8 & 随机量子化的新生：AI 采样 & 3\\
\bottomrule
\end{tabular}
\end{center}

贯穿全书并在第六章会合的有两条主线：
\begin{itemize}
  \item \textbf{梯度流主线}（L\"uscher 纲领）：光滑规范场作为重整化场，
        从平庸化映射到涂抹准分布；
  \item \textbf{大动量主线}（Ji 的 LaMET 与 Radyushkin 的伪分布）：
        高速等时关联函数经微扰匹配得到光锥分布。
\end{itemize}
环绕它们的是使能技术（第二章）、让第一性原理 PDF 成为可能的硬核
重整化问题（第五章），以及格点结果通往物理世界的两个接口：全球 PDF
拟合（第七章）与植根于随机量子化的机器学习采样器（第八章）。

\section*{如何阅读各条目}

每条条目以一张书目卡片开篇，随后回答三个问题：\emph{为何}入选
（附其在本语料库内的枢纽地位）、\emph{究竟说了什么}（物理内容与关键
公式），以及它的\emph{回响}出现在库内何处。条目自成一体；公式按现代
记法转写以便定位，而非原文复刻。

\bigskip
\noindent\emph{诚实声明。} 下文的引用计数均指在本 99 篇语料库内
（``$\sim$N/99''）的计数，取保守值。书目数据转录自库内论文自身的参考
文献列表并与标准记录交叉核对；无法确认的字段标注 [?]。
